\section{Docker}
Docker \cite{Docker} jest narzędziem do konteneryzacji, która umożliwia tworzenie i uruchamianie aplikacji w kontenerach, które są odizolowanymi środowiskami. Kontenery są zminimalizowanymi obrazami, co przyczynia się do ich lekkości i reużywalności. Zawierają tylko potrzebne rzeczy do prawidłowego działania aplikacji. 
\subsection{Jak działa Docker?}
Docker polega na zasadzie wirtualizacji na poziomie systemu operacyjnego. Poniżej wypunktowano główne narzędzia dockera:
\begin{enumerate}
    \item \textbf{Obrazy}: są to szablony zawierające system operacyjny i podstawowe zależności do działania aplikacji.
    \item \textbf{Kontenery}:
    obraz może zostać użyty i uruchomiony jako kontener, który jest odizolowanym środowiskiem, dzięki czemu każdy może pracować na tym samym środowisku. Jest to przydatne w celu wyeliminowania stawiania i konfigurowania środowisk dla każdej osoby. W dodatku przez różne systemy operacyjne aplikacja mogłaby działać w inny sposób, docker eleminuje ten problem.
\end{enumerate}

\subsection{Zalety Dockera}
\begin{enumerate}
    \item \textbf{Przenośność}: Kontenery mogą być uruchamiane na różnych systemach operacyjnych i infrastrukturach, co ułatwia wdrażanie aplikacji.
    Docker może zostać użyty na głównych systemach operacyjnych, co znacząco ułatwia wdrażanie aplikacji. Mac i linux posiadają paczki do instalacji dockera, natomiast na windowsie można użyć aplikacji docker deskop i dodatkowo zainstalować podsystem linux WSL2 (Windows subsystem for linux), w celu zwiększenia wydajności kontenerów.
    \item \textbf{Izolacja}:
    Aplikacja działa w odizolowanym środowisku, co minimalizuje konflikty między środowiskami i zależnościami.
    \item \textbf{Szybkość}: Kontenery są lżejsze i szybsze do uruchomienia w porównaniu do tradycyjnych maszyn wirtualnych.
    Kontenery są lekkie i szybkie do uruchomienia w porównaniu do maszyn wirtualnych. Kontenery używają tylko najpotrzebniejszych zależności, co przyczynia się do ich lekkości.
    \item \textbf{Skalowalność}: Docker umożliwia w prosty sposób skalowanie aplikacji, poprzez uruchamianie wielu kontenerów jednocześnie. Np. na jednym kontenerze mamy aplikację webową, a na drugim bazę danych.
\end{enumerate}
\subsection{Obraz i kontener użyty w aplikacji}
Przy budowaniu obrazu został rozszerzony obraz \textit{node:22-alpine}, \textit{alpine} oznacza okrojoną wersję, która posiada tylko najpotrzebniejsze zależności. Na powyższym obrazie jest dostępny Node.js w wersji 22. Dodatkowo został zainstalowany bash, który ułatwia poruszanie się po zbudowanym kontenerze. Do zbudowania obrazu zostały użyte polecenia instalujące paczki npm potrzebne do prawidłowego działania aplikacji, dodatkowo instalowana jest globalnie paczka serve, która służy do obsługi serwera HTTP. Po włączeniu kontenera inicjowane są polecenia uruchamiające websocket na porcie \textbf{8080} jak i aplikacje webową, przy użyciu serwera HTTP, na porcie \textbf{3000}. Powyższe porty są udostępniane dla kontenera poleceniem \textit{EXPOSE}. W pliku konfiguracyjnym dla kontenera zostały zdefiniowane porty poza kontenerem, które są widoczne przez komputer. Są takie same jak w kontenerze. Nazwa kontenera została ustawiona na app\_container. Dodatkowo została udostępniona sieć dla kontenera, w celu prawidłowego połączenia się z serwerem na STM32.
\subsection{Uruchomienie aplikacji w kontenerze}
Całą aplikację można uruchomić poleceniem \textit{docker compose up}. Po użyciu powyższej komendy na porcie \textbf{3000} jest dostępna aplikacja z intefejsem do obsługi teleskopu i WebSocket na porcie \textbf{8080}. W celu wejścia na kontener należy użyć polecenia \textit{docker exec -it app\_container bash}.